\section{Muestreo aleatorio}

\subsection{Motivación: Estimación de parámetros poblacionales}

Supongamos que tratamos de encontrar la edad promedio en una ciudad, digamos Oaxaca. Una manera de hacerlo sería por \emph{fuerza bruta}, es decir, recolectando esta información persona por persona. Pero este método sería muy costoso en términos de infraestructura y tiempo.

En estadística, este es un problema común, cuya solución está en el \emph{muestreo aleatorio}: Tomemos un grupo de 1000 individuos (o 10,000 dependiendo de tu capacidad, obviamente entre más, es mejor) y calculemos la edad promedio en este grupo, a la que denotaremos por $A_{1}.$

Repitamos este procedimiento, digamos 100 veces, y denotaremos por $A_{1}, A_{2},...,A_{100}$ el promedio de edades obtenido en cada respectivo intento.

\subsection{Ley de los grandes números}

De acuerdo a la \emph{ley de los grandes números}, la cantidad
\begin{align}
	\bar{A}_{100}=\dfrac{A_{1}+...+A_{100}}{100}
\end{align}
es una aproximación muy cercana al promedio real de la edad de los pobladores de la ciudad.

\begin{observacion}
	No estamos más interesados en obtener el valor exacto de la edad promedio, sino establecer un \emph{estimador} para la misma. En tal caso, tenemos que conformarnos con la definición de un \emph{rango de valores} en el que el valor real podría estar.
\end{observacion}

\subsection{Teorema del límite central}

De acuerdo al \emph{teorema del límite central}, si el número de tales muestras es suficientemente grande, $A_{1},A_{2},...,A_{100}$ estarán distribuidos de manera normal.

\begin{teorema}[Teorema del límite central]
	Si $X_1, X_2, \ldots, X_n$ son variables aleatorias independientes e idénticamente distribuidas con media $\mu$ y varianza $\sigma^2$, entonces cuando $n$ es suficientemente grande:
	\begin{align}
		\frac{\bar{X} - \mu}{\sigma/\sqrt{n}} \xrightarrow{d} N(0,1)
	\end{align}
	
	Es decir, la distribución de la media muestral estandarizada se aproxima a una distribución normal estándar.
\end{teorema}

\subsection{Implicaciones prácticas del TLC}

El teorema del límite central tiene profundas implicaciones para la inferencia estadística:

\begin{enumerate}
	\item \textbf{Independencia de la distribución original:} No importa si la población original sigue una distribución normal, exponencial, uniforme o cualquier otra. La distribución de la media muestral será aproximadamente normal para muestras grandes.
	
	\item \textbf{Tamaño de muestra:} En la práctica, $n \geq 30$ suele ser suficiente para que la aproximación normal sea buena, aunque esto depende de la forma de la distribución original.
	
	\item \textbf{Parámetros de la distribución muestral:}
	\begin{align}
		E(\bar{X}) &= \mu \\
		\text{Var}(\bar{X}) &= \frac{\sigma^2}{n} \\
		\text{DE}(\bar{X}) &= \frac{\sigma}{\sqrt{n}}
	\end{align}
\end{enumerate}

\subsection{Ejemplo ilustrativo}

Consideremos el lanzamiento de un dado justo. La distribución de un solo lanzamiento es uniforme discreta en $\{1, 2, 3, 4, 5, 6\}$ con:
\begin{align}
	\mu &= \frac{1+2+3+4+5+6}{6} = 3.5 \\
	\sigma^2 &= \frac{(1-3.5)^2 + (2-3.5)^2 + \cdots + (6-3.5)^2}{6} = \frac{35}{12} \approx 2.917
\end{align}

Si lanzamos el dado $n = 36$ veces y calculamos el promedio $\bar{X}$, por el TLC:
\begin{align}
	\bar{X} \sim N\left(3.5, \frac{2.917}{36}\right) = N(3.5, 0.081)
\end{align}

Esto significa que:
\begin{align}
	P(3.0 < \bar{X} < 4.0) &= P\left(\frac{3.0-3.5}{\sqrt{0.081}} < Z < \frac{4.0-3.5}{\sqrt{0.081}}\right) \\
	&= P(-1.76 < Z < 1.76) \\
	&\approx 0.921
\end{align}

Hay aproximadamente un 92.1\% de probabilidad de que el promedio de 36 lanzamientos esté entre 3.0 y 4.0.

\subsection{Simulación del TLC con Python}

Podemos verificar el teorema del límite central mediante simulación:

\begin{lstlisting}[language=Python]
import numpy as np
import matplotlib.pyplot as plt

# Parámetros
n = 30  # tamaño de muestra
num_simulaciones = 10000

# Simulación: promedio de n valores de una distribución exponencial
medias_muestrales = []
for _ in range(num_simulaciones):
    muestra = np.random.exponential(scale=2.0, size=n)
    medias_muestrales.append(np.mean(muestra))

# Graficar histograma
plt.hist(medias_muestrales, bins=50, density=True, alpha=0.7)

# Superponer curva normal teórica
mu = 2.0  # media de la exponencial
sigma = 2.0  # desviación estándar de la exponencial
x = np.linspace(mu - 4*sigma/np.sqrt(n), mu + 4*sigma/np.sqrt(n), 100)
y = (1/(sigma/np.sqrt(n) * np.sqrt(2*np.pi))) * np.exp(-0.5*((x-mu)/(sigma/np.sqrt(n)))**2)
plt.plot(x, y, 'r-', linewidth=2)

plt.xlabel('Media muestral')
plt.ylabel('Densidad')
plt.title(f'Distribución de medias muestrales (n={n})')
plt.show()
\end{lstlisting}

\subsection{Aplicaciones del TLC}

El teorema del límite central es fundamental en muchas áreas:

\begin{itemize}
	\item \textbf{Control de calidad:} Las medias de mediciones de calidad siguen distribuciones normales, permitiendo el uso de gráficos de control.
	
	\item \textbf{Encuestas y sondeos:} Las proporciones muestrales (que son promedios de variables binarias) siguen distribuciones normales aproximadas.
	
	\item \textbf{Finanzas:} Los rendimientos promedio de carteras de inversión se modelan como normales.
	
	\item \textbf{Ciencias naturales:} Muchas mediciones biológicas y físicas son el resultado de muchos factores pequeños e independientes, por lo que siguen distribuciones normales.
\end{itemize}

\subsection{Limitaciones y consideraciones}

Aunque el TLC es muy poderoso, es importante recordar:

\begin{enumerate}
	\item \textbf{Tamaño de muestra:} Para distribuciones muy asimétricas o con colas pesadas, puede requerirse $n > 30$ para una buena aproximación.
	
	\item \textbf{Independencia:} Las observaciones deben ser independientes. Si hay correlación entre observaciones, el TLC puede no aplicarse.
	
	\item \textbf{Varianza finita:} El TLC requiere que la varianza poblacional sea finita. Para distribuciones con varianza infinita (como la distribución de Cauchy), el TLC no se aplica.
\end{enumerate}

\subsection{Resumen}

El muestreo aleatorio y el teorema del límite central son la columna vertebral de la inferencia estadística:

\begin{itemize}
	\item El \textbf{muestreo aleatorio} nos permite hacer inferencias sobre poblaciones grandes a partir de muestras pequeñas.
	
	\item La \textbf{ley de los grandes números} garantiza que nuestras estimaciones convergen a los valores verdaderos.
	
	\item El \textbf{teorema del límite central} nos permite usar la distribución normal para hacer inferencias, incluso cuando la población original no es normal.
\end{itemize}

Estos resultados teóricos justifican las técnicas de estimación y pruebas de hipótesis que estudiaremos en las siguientes secciones.
